% ============================================================================
% Chapter 12: Distributed Training and Multi-Node Inference
% ============================================================================

\chapter{Distributed Training and Multi-Node Inference}
\label{ch:distributed}

\begin{quote}
\textit{``Scalability is not about adding more hardware---it's about adding more information per byte transferred.''}
\end{quote}

\section{The Bandwidth Bottleneck in Distributed Training}
\label{sec:dist_bandwidth}

In distributed deep learning, the dominant cost is not computation but communication. For data-parallel training with $P$ workers, each synchronization step requires broadcasting the complete gradient vector $\nabla \ell$ of dimension $d$. For FP32:

\[
B_{\text{sync}} = 2 \cdot d \cdot P \text{ bytes per all-reduce}
\]

For a 64M-parameter model with $P = 4$ workers:

\[
B_{\text{sync}} = 2 \times 64 \times 10^6 \times 4 = 512 \text{ MB per step}
\]

At 10\,Gbps inter-node bandwidth (typical for 10GbE), this takes 410\,ms---completely dominating the compute time of approximately 50\,ms per step on consumer hardware.

\section{OIL-Optimized Gradient Compression}
\label{sec:oil_gradient_compression}

MYTHOS addresses this bottleneck by transmitting gradients in OIL format. Since the codebook structure is shared across all workers (fixed format per layer), only the codebook and packed indices need to be transmitted:

\[
B_{\text{sync}}^{\text{OIL}} = d \times \frac{\BPW}{8} + K \times 4 + \text{header}
\]

For OIL4 ($\BPW = 4$, $K = 16$):

\[
B_{\text{sync}}^{\text{OIL4}} = 64 \times 10^6 \times 0.5 + 16 \times 4 + 64 = 32.00 \text{ MB}
\]

This represents a 16$\times$ reduction in synchronization bandwidth compared to FP32. The communication time drops from 410\,ms to 25.6\,ms, making the all-reduce negligible compared to the compute time.

\section{Topology-Aware Communication}
\label{sec:topology}

\subsection{Ring All-Reduce with Format Awareness}

The standard ring all-reduce algorithm is modified to be format-aware:

\begin{algorithm}[H]
\caption{Format-Aware Ring All-Reduce}
\label{alg:ring_allreduce}
\begin{algorithmic}[1]
\STATE \textbf{Input:} Local gradient $g_i$ at worker $i$, format descriptor $F$
\STATE Pack $g_i$ into OIL format $g_i^{\text{OIL}}$ using format $F$
\STATE $S \leftarrow |g_i^{\text{OIL}}|$ (packed size)
\FOR{step $= 1$ to $P-1$}
    \STATE Send chunk $(\text{step} \bmod P) \cdot S/P$ to next ring node
    \STATE Receive chunk from previous ring node
    \STATE Accumulate: $g_{\text{recv}} \to g_{\text{acc}}$
\ENDFOR
\STATE Unpack $g_{\text{acc}}$ from OIL format
\STATE \textbf{Output:} Synchronized gradient $g_{\text{acc}}$
\end{algorithmic}
\end{algorithm}

The ring all-reduce requires $2(P-1)/P$ messages of size $d \cdot \BPW / (8P)$ each. For $P = 4$ and OIL4, each message is 8\,MB---well within the optimal message size for TCP-based networks.

\subsection{Hierarchical All-Reduce}

For multi-node setups ($P > 8$), a two-level hierarchical all-reduce is used:

\begin{enumerate}
    \item \textbf{Intra-node}: GPU-to-GPU via PCIe P2P (or NVLink on supported hardware) at full bandwidth.
    \item \textbf{Inter-node}: CPU-mediated all-reduce over Ethernet with OIL-compressed gradients.
\end{enumerate}

The hierarchical approach reduces the inter-node communication volume by a factor of $N_{\text{gpus\_per\_node}}$, since intra-node communication is approximately 10$\times$ faster than inter-node.

\section{Gradient Quantization for Communication}
\label{sec:grad_quant}

\subsection{The Double-Quantization Insight}

In distributed MYTHOS training, gradients undergo a unique double-quantization:

\begin{enumerate}
    \item \textbf{Local quantization}: Each worker quantizes its local gradient using the shared codebook.
    \item \textbf{Communication}: The quantized (packed) gradients are transmitted.
    \item \textbf{Dequantization}: Each worker dequantizes the aggregated gradient.
\end{enumerate}

The critical insight is that gradient quantization error is \textbf{identical across workers} because they share the same codebook. This means the quantization bias cancels out during the all-reduce:

\[
\frac{1}{P} \sum_{i=1}^P Q(g_i) = Q\left(\frac{1}{P} \sum_{i=1}^P g_i\right) + \frac{1}{P} \sum_{i=1}^P \epsilon_i
\]

where $\epsilon_i = Q(g_i) - g_i$ is the quantization error. Since $\E[\epsilon_i] = 0$ (Lloyd-Max property), the expected aggregation error is $\sigma^2_\epsilon / P$, which vanishes with more workers.

\begin{mythm}[Communication-Efficient Quantized Aggregation]{thm:comm_quant}
Let $g_1, \ldots, g_P$ be gradients at $P$ workers sharing a Lloyd-Max codebook $\mathcal{C}$. The quantized all-reduce $\bar{g} = \frac{1}{P} \sum_{i=1}^P Q_{\mathcal{C}}(g_i)$ satisfies $\|\bar{g} - \bar{g}_{\text{true}}\|_2^2 \leq \sigma^2_{\text{QMSE}} / P$, where $\sigma^2_{\text{QMSE}}$ is the single-worker quantization MSE and $\bar{g}_{\text{true}} = \frac{1}{P} \sum_{i=1}^P g_i$.
\end{mythm}

\begin{proof}
The quantization error at worker $i$ is $\epsilon_i = Q(g_i) - g_i$. By the Lloyd-Max optimality condition, $\E[\epsilon_i] = 0$ and $\Var[\epsilon_i] = \sigma^2_{\text{QMSE}}$. The aggregation error is:

\[
\bar{g} - \bar{g}_{\text{true}} = \frac{1}{P}\sum_{i=1}^P \epsilon_i
\]

Since the $\epsilon_i$ are approximately independent across workers (different data samples):

\[
\left\|\frac{1}{P}\sum_{i=1}^P \epsilon_i\right\|_2^2 = \frac{1}{P^2} \cdot P \cdot \sigma^2_{\text{QMSE}} = \frac{\sigma^2_{\text{QMSE}}}{P}
\]

$\square$
\end{proof}

\section{Fault Tolerance}
\label{sec:dist_fault}

\subsection{Weight Checkpointing with OIL Snapshots}

The OIL format's compact representation enables frequent, low-cost weight checkpoints:

\begin{longtable}{lcr}
\toprule
\textbf{Checkpoint Format} & \textbf{Size (64M)} & \textbf{Write Time} \\
\midrule
\endhead
FP32 raw & 256\,MB & 67\,ms \\
OIL4 snapshot & 32\,MB & 8\,ms \\
OIL2 snapshot & 16\,MB & 4\,ms \\
\bottomrule
\end{longtable}

At OIL4, checkpoints can be written every 100 training steps with negligible overhead (8\,ms vs.\ 50\,ms per step). This provides automatic recovery from any single-worker failure with at most 100 steps of lost progress.

\subsection{Elastic Scaling}

When workers join or leave, MYTHOS performs elastic scaling:

\begin{enumerate}
    \item Paused training at current step.
    \item Redistribute gradient chunks across surviving workers.
    \item Adjust learning rate using linear scaling rule: $\eta_{\text{new}} = \eta_0 \cdot P_{\text{new}} / P_{\text{old}}$.
    \item Resume training from last checkpoint.
\end{enumerate}

The OIL format's fixed codebook structure means no re-initialization is required: the new workers simply adopt the shared codebook and begin processing their assigned gradient chunks.

\section{Multi-Node Inference}
\label{sec:multi_node_inference}

For models too large for a single node's memory, MYTHOS supports tensor-parallel inference:

\begin{enumerate}
    \item Split weight matrix $W \in \mathbb{R}^{M \times N}$ along columns: $W = [W_1 \mid W_2 \mid \cdots \mid W_P]$.
    \item Each node $i$ holds $W_i$ in OIL format.
    \item For input $x$, each node computes the partial sum $y_i = W_i x_i$ where $x_i$ is the corresponding slice of $x$.
    \item An all-gather assembles the complete output $y = [y_1 \mid y_2 \mid \cdots \mid y_P]$.
\end{enumerate}

The communication volume per inference step is $|x| \times P$ bytes (input broadcast) plus $|y| \times P$ bytes (output gather). For a 64M model with hidden size 1024:

\[
B_{\text{infer}} = 2 \times 1024 \times 4 \times P = 8192P \text{ bytes}
\]

At $P = 4$, this is 32\,KB per token---negligible even over 1\,Gbps Ethernet.

\section{Distributed Benchmarks}
\label{sec:dist_benchmarks}

\begin{longtable}{llccc}
\toprule
\textbf{Config} & \textbf{Format} & \textbf{Steps/s} & \textbf{Sync ms} & \textbf{Efficiency} \\
\midrule
\endhead
1 node, 1 GPU & OIL4 & 18.2 & 0 & 100\% \\
2 node, 2 GPU & OIL4 & 16.8 & 12 & 92\% \\
4 node, 4 GPU & OIL4 & 15.1 & 26 & 83\% \\
4 node, 4 GPU & FP32 & 8.4 & 410 & 46\% \\
8 node, 8 GPU & OIL4 & 13.2 & 28 & 73\% \\
8 node, 8 GPU & OIL2 & 14.1 & 14 & 77\% \\
\bottomrule
\end{longtable}

The OIL4 distributed training achieves 83\% parallel efficiency at 4 nodes, compared to only 46\% for FP32---nearly double the effective throughput at the same hardware cost.

\section{Network Requirements}
\label{sec:network_reqs}

For optimal distributed MYTHOS training:

\begin{longtable}{lc}
\toprule
\textbf{Topology} & \textbf{Max Workers} \\
\midrule
Loopback (single node) & 4 (CPU threads) \\
1\,Gbps Ethernet & 8 \\
10\,Gbps Ethernet & 32 \\
InfiniBand HDR & 256+ \\
\bottomrule
\end{longtable}

The OIL format's bandwidth efficiency means that even consumer-grade networking equipment can support meaningful multi-node training, democratizing distributed AI development.
